% 版式
\usepackage{geometry}
\geometry{
  a4paper,
  margin=2.5cm,
  headheight=14pt,
}

% 主题色（克制、印刷友好）
\usepackage{xcolor}
\definecolor{themeaccent}{HTML}{1a5f7a}
\definecolor{themecolorthm}{HTML}{1a5f7a}
\definecolor{themecolordef}{HTML}{2d6a4f}
\definecolor{themecolorremark}{HTML}{8b6914}

% 定理块底纹浓度（xcolor：数字越大底色越深；旧版约 !4–!6，此处加深便于对比）
\newcommand{\thmshadecolor}{themecolorthm!18}
\newcommand{\defshadecolor}{themecolordef!14}
\newcommand{\remarkshadecolor}{themecolorremark!12}

% 数学
\usepackage{amsmath,amssymb}
\usepackage{mathtools}
\usepackage{amsthm}
\usepackage[anycase]{thmtools}
\usepackage{bm}

% 行距（可读性）
\usepackage{setspace}
\onehalfspacing

% 页眉页脚
\usepackage{fancyhdr}
\pagestyle{fancy}
\fancyhf{}
\fancyhead[L]{\small\color{themeaccent!70!black}\nouppercase{\leftmark}}
\fancyhead[R]{\small\thepage}
\renewcommand{\headrulewidth}{0.4pt}
\renewcommand{\headrule}{\color{themeaccent!50!black}\hrule height 0.4pt}

% 章节标题着色 + 装饰线（ctex）
\ctexset{
  section/format     = \Large\bfseries\color{themeaccent},
  section/aftername  = \hspace{0.5em},
  section/aftertitle = \par\vspace{-4pt}%
    {\color{themeaccent!40}\rule{\linewidth}{0.5pt}}\vspace{4pt},
  subsection/format  = \large\bfseries\color{themeaccent!88!black},
}

% 常用配对符号（mathtools）
\DeclarePairedDelimiter{\paren}{(}{)}
\DeclarePairedDelimiter{\sqbrk}{[}{]} % 勿用 \brack，与内核冲突
\DeclarePairedDelimiter{\curlbrace}{\{}{\}} % 勿用 \brace，与内核冲突
\DeclarePairedDelimiter{\abs}{\lvert}{\rvert}
\DeclarePairedDelimiter{\norm}{\lVert}{\rVert}

% 定理样式：彩色标题（底纹在 \declaretheorem 中用 shaded 指定）
\declaretheoremstyle[
  spaceabove=6pt,
  spacebelow=6pt,
  headfont=\bfseries\color{themecolorthm},
  notefont=\bfseries\color{themecolorthm!80!black},
  bodyfont=\itshape,
  headpunct={.},
  postheadspace=0.5em,
]{styletheorem}

\declaretheoremstyle[
  spaceabove=6pt,
  spacebelow=6pt,
  headfont=\bfseries\color{themecolordef},
  notefont=\bfseries\color{themecolordef!80!black},
  bodyfont=\normalfont,
  headpunct={.},
  postheadspace=0.5em,
]{styledefinition}

\declaretheoremstyle[
  spaceabove=6pt,
  spacebelow=6pt,
  headfont=\bfseries\color{themecolorremark},
  notefont=\bfseries\color{themecolorremark!80!black},
  bodyfont=\normalfont,
  headpunct={.},
  postheadspace=0.5em,
]{styleremark}

\declaretheorem[
  style=styletheorem,
  name=定理,
  numberwithin=section,
  shaded={bgcolor=\thmshadecolor, margin=0pt},
]{theorem}

\declaretheorem[
  style=styletheorem,
  name=引理,
  sibling=theorem,
  shaded={bgcolor=\thmshadecolor, margin=0pt},
]{lemma}

\declaretheorem[
  style=styletheorem,
  name=命题,
  sibling=theorem,
  shaded={bgcolor=\thmshadecolor, margin=0pt},
]{proposition}

\declaretheorem[
  style=styletheorem,
  name=推论,
  sibling=theorem,
  shaded={bgcolor=\thmshadecolor, margin=0pt},
]{corollary}

\declaretheorem[
  style=styledefinition,
  name=定义,
  sibling=theorem,
  shaded={bgcolor=\defshadecolor, margin=0pt},
]{definition}

\declaretheorem[
  style=styleremark,
  name=注,
  sibling=theorem,
  shaded={bgcolor=\remarkshadecolor, margin=0pt},
]{remark}

\renewcommand{\proofname}{\textcolor{themecolorthm}{证明}}
\renewcommand{\qedsymbol}{\textcolor{themecolorthm}{$\blacksquare$}}

% 彩色标题页（替代 \maketitle）
\makeatletter
\newcommand{\maketitlecolored}{%
  \begin{center}\leavevmode
    {\LARGE\bfseries\color{themeaccent}\@title\par}
    \vspace{0.8em}
    {\large\@author\par}
    {\normalsize\color{black!55}\@date\par}
  \end{center}
  \vspace{1.8em}%
}
\makeatother

% 表格（专业三线表）
\usepackage{booktabs}

% 图表标题样式
\usepackage{caption}
\captionsetup{labelfont={bf,color=themeaccent}, labelsep=period, font=small}

% 链接（置于 preamble 末尾）
\usepackage{hyperref}
\hypersetup{
  colorlinks=true,
  linkcolor=themeaccent!90!black,
  citecolor=themecolordef!75!black,
  urlcolor=themeaccent,
}

% 智能交叉引用（须在 hyperref 之后加载）
\usepackage{cleveref}
\crefname{theorem}{定理}{定理}
\crefname{lemma}{引理}{引理}
\crefname{proposition}{命题}{命题}
\crefname{corollary}{推论}{推论}
\crefname{definition}{定义}{定义}
\crefname{remark}{注}{注}
\crefformat{equation}{式~(#2#1#3)}

% ============================================================
% 以下为本文（verl 笔记）额外追加的部分，上面的模板内容未作改动
% ============================================================

% 插图
\usepackage{graphicx}
\graphicspath{{figures/}}

% ------------------------------------------------------------
% 深层编号列表：原笔记有 4–6 层 tab 缩进。
% enumitem 内置的 enumerate 被 \thr@@(=4) 限制，单靠 \setlistdepth 不够，
% 故用 \newlist 定义一个 max-depth=9 的 enum 列表环境。
% ------------------------------------------------------------
\usepackage{enumitem}
\setlistdepth{9}
% 用 \renewlist 重定义标准的 enumerate / itemize，把深度上限从 4 提到 9，
% 这样正文里原有的 \begin{enumerate} 无需逐个改名。
\renewlist{enumerate}{enumerate}{9}
\setlist[enumerate,1]{label=\arabic*., leftmargin=2em, itemsep=2pt, topsep=4pt}
\setlist[enumerate,2]{label=(\arabic*), leftmargin=2em, itemsep=2pt, topsep=3pt}
\setlist[enumerate,3]{label=\roman*., leftmargin=2em, itemsep=1pt, topsep=2pt}
\setlist[enumerate,4]{label=\alph*., leftmargin=2em, itemsep=1pt, topsep=2pt}
\setlist[enumerate,5]{label=\textbullet, leftmargin=1.6em, itemsep=1pt, topsep=2pt}
\setlist[enumerate,6]{label=--, leftmargin=1.6em, itemsep=1pt, topsep=2pt}
\setlist[enumerate,7]{label=$\cdot$, leftmargin=1.6em, itemsep=1pt, topsep=2pt}
\setlist[enumerate,8]{label=$\ast$, leftmargin=1.6em, itemsep=1pt, topsep=2pt}
\setlist[enumerate,9]{label=$\circ$, leftmargin=1.6em, itemsep=1pt, topsep=2pt}
\renewlist{itemize}{itemize}{9}
\setlist[itemize,1]{label=\textbullet, leftmargin=2em, itemsep=2pt, topsep=3pt}
\setlist[itemize,2]{label=--, leftmargin=2em, itemsep=1pt, topsep=2pt}
\setlist[itemize,3]{label=$\cdot$, leftmargin=1.6em, itemsep=1pt, topsep=2pt}

% ------------------------------------------------------------
% 行内代码：\detokenize 把 _ % & # 等一律转成「其它」字符再排版，
% 不会触发数学下标或注释。注意参数里不要放中文（中文移到 \code 外面）；
% 也不要放 Unicode 数学符号（等宽字体缺字形，改用 LaTeX 命令或 ASCII）。
% ------------------------------------------------------------
\newcommand{\code}[1]{{\ttfamily\detokenize{#1}}}

% 新增内容标记：灰字 + 【补充】标签，与原文在视觉上可区分
% 用 \color 分组而不是 \textcolor{...}{...}：后者把段落放进参数里，
% make4ht/tex4ht 会报 "Argument of \c:HColor: has an extra }"。
\newcommand{\note}[1]{\par\smallskip\noindent
  \textbf{【补充】}~#1\par\smallskip}
% 行内新增内容标记（不另起段）
\newcommand{\inlinenote}[1]{〔补充：#1〕}
% 表格中换行
\usepackage{array}
\newcolumntype{L}[1]{>{\raggedright\arraybackslash}p{#1}}

% 符号简写（\mathbb 由 amssymb 提供）
\newcommand{\KL}{\mathrm{KL}}
\newcommand{\GAE}{\mathrm{GAE}}
\newcommand{\R}{\mathbb{E}}
% Unicode 圈号：Latin 字体没有，用带圈数字代替
\newcommand{\stagenum}[1]{\textcircled{\scriptsize #1}}
